\chapter{Fondamenti teorici}
\label{cap:fondamenti}

\todoscrivi{Le basi necessarie a leggere il resto della tesi. Non un manuale:
includere solo cio' che verra' effettivamente usato nei capitoli successivi.}

\section{Reti neurali generative}
\label{sec:reti-generative}

\section{Generative Adversarial Networks}
\label{sec:gan}
% Formulazione minimax, ruolo di generatore e discriminatore, criticita' note
% dell'addestramento (mode collapse, instabilita').

\section{Prospettive teoriche sulla creativita'}
\label{sec:teoria-creativita}
% Creativita' combinatoria, esplorativa, trasformazionale (Boden).
% Criteri/tripode della creativita' computazionale (Colton).
% Distinzione tra creativita' e percezione della creativita'.
% Base concettuale necessaria PRIMA di presentare la rivendicazione delle CAN
% (sezione successiva): senza un criterio di cosa conterebbe come creativita',
% la sezione sul caso CAN non avrebbe contro cosa misurarsi.

\section{Creative Adversarial Networks: il caso di partenza}
\label{sec:can}
% ADR-0007 (2026-09-14): presentata come CASO discusso dalla letteratura
% (Elgammal et al., 2017), non riprodotta sperimentalmente da Gian in questo
% documento. La rivendicazione di creativita' non e' implicita: sta nel nome
% dell'architettura e nella funzione di perdita (loss di ambiguita' stilistica).
% E' il punto di partenza scelto proprio per questa esplicitezza -- apre il
% capitolo 3 (evoluzione dell'argomentazione critica).
